% Capítulo 3: Metodología
% Propósito: Detallar el enfoque metodológico, diseño de investigación, técnicas e instrumentos de recolección de información y la metodología de desarrollo de software seleccionada.
\section{Marco Metodológico}

\section{Cronograma de actividades}
El siguiente cronograma organiza el desarrollo del sistema durante un período de 13 semanas, equivalentes a 65 días hábiles. Las primeras dos semanas estarán destinadas al levantamiento de requerimientos, planificación y preparación técnica del proyecto. Posteriormente, la ejecución de las actividades se gestionará mediante Kanban, permitiendo que cada funcionalidad avance de forma continua a través de los estados definidos en el tablero, respetando los límites de trabajo en progreso y los criterios de validación antes de considerarse completada.

Fase 1: Planificación y preparación del proyecto (Semanas 1–2)
Semana 1: Levantamiento y análisis de requerimientos.
Se realizará el diagnóstico inicial del proceso de gestión de incidencias en la Alcaldía Municipal de Rivas, identificando las necesidades de los ciudadanos y del personal encargado de administrar los reportes. A partir de esta información se definirán y documentarán los requerimientos funcionales y no funcionales de la plataforma.

Semana 2: Planificación técnica y organización del trabajo.
Se priorizarán los requerimientos identificados y se descompondrán en tareas para su gestión mediante Kanban. Durante esta semana se configurará el tablero de trabajo, se establecerán los estados de las tareas y los límites de trabajo en progreso. Asimismo, se definirá la arquitectura inicial del sistema, el modelo de datos y se prepararán los entornos de desarrollo utilizando Angular, Node.js y PostgreSQL.

Fase 2: Desarrollo de funcionalidades base (Semanas 3–6)
Semanas 3–4: Gestión de usuarios, roles y autenticación.
Se desarrollarán las funcionalidades relacionadas con registro e inicio de sesión, gestión de roles y permisos, autenticación mediante tokens y validaciones de seguridad. El trabajo incluirá backend, frontend y pruebas de las funcionalidades desarrolladas.

Semanas 5–6: Registro de incidencias comunitarias.
Se desarrollará el módulo de registro de incidencias, incluyendo descripción, categoría, ubicación geográfica, evidencia fotográfica y las validaciones necesarias para garantizar la integridad de la información.

Fase 3: Funcionalidades ciudadanas y administrativas (Semanas 7–9)
Semana 7: Seguimiento de incidencias por parte del ciudadano.
Se implementará la consulta del historial de reportes, visualización detallada y seguimiento del estado de cada incidencia.

Semanas 8–9: Gestión administrativa de incidencias.
Se desarrollará el panel destinado al personal municipal, incluyendo consulta, filtrado, asignación, actualización de estados e historial de cambios de los reportes.

Fase 4: Visualización, validación y cierre (Semanas 10–13)
Semanas 10–11: Visualización territorial y métricas.
Se implementarán los mapas para representar incidencias geográficamente, mecanismos para identificar concentraciones de reportes y métricas básicas para apoyar la gestión municipal.

Semana 12: Pruebas integrales y validación.
Se realizarán pruebas de integración, seguridad, usabilidad y funcionamiento general del sistema, verificando los principales flujos de uso y corrigiendo los problemas detectados.

Semana 13: Optimización, documentación y entrega.
Se realizarán los ajustes finales, documentación técnica, manuales de usuario, preparación del despliegue y presentación de la versión final.

